% Backtesting Engine: Complete Educational Guide
\documentclass[11pt,a4paper]{article}

% Packages
\usepackage[T1]{fontenc}
\usepackage[utf8]{inputenc}
\usepackage{lmodern}
\usepackage{geometry}
\usepackage{hyperref}
\usepackage{titlesec}
\usepackage{enumitem}
\usepackage{xcolor}
\usepackage{graphicx}
\usepackage{amsmath,amssymb}
\usepackage{listings}
\usepackage{caption}
\usepackage{tcolorbox}
\usepackage{multicol}
\usepackage{fancyhdr}

\geometry{margin=1in}
\hypersetup{
  colorlinks=true,
  linkcolor=blue,
  urlcolor=blue,
  citecolor=blue
}

% Custom boxes for key concepts
\newtcolorbox{keyterm}[1]{
  colback=blue!5!white,
  colframe=blue!75!black,
  fonttitle=\bfseries,
  title=#1
}

\newtcolorbox{example}[1]{
  colback=green!5!white,
  colframe=green!75!black,
  fonttitle=\bfseries,
  title=#1
}

\newtcolorbox{warning}{
  colback=red!5!white,
  colframe=red!75!black,
  fonttitle=\bfseries,
  title=Important Note
}

\newtcolorbox{flowstep}[1]{
  colback=yellow!5!white,
  colframe=orange!75!black,
  fonttitle=\bfseries,
  title=Step #1
}

% Code listing style
\lstdefinestyle{pythoncode}{
  language=Python,
  basicstyle=\ttfamily\small,
  keywordstyle=\color{blue}\bfseries,
  commentstyle=\color{gray}\itshape,
  stringstyle=\color{red},
  breaklines=true,
  showstringspaces=false,
  frame=single,
  numbers=left,
  numberstyle=\tiny\color{gray},
  rulecolor=\color{black!20}
}

\lstdefinestyle{jsoncode}{
  basicstyle=\ttfamily\small,
  breaklines=true,
  showstringspaces=false,
  frame=single,
  numbers=left,
  numberstyle=\tiny\color{gray},
  rulecolor=\color{black!20}
}

\lstdefinestyle{shellcode}{
  language=bash,
  basicstyle=\ttfamily\small,
  breaklines=true,
  showstringspaces=false,
  frame=single,
  rulecolor=\color{black!20}
}

\title{%
  \textbf{Backtesting Engine}\\
  \large Complete Study Guide with Key Concepts, Architecture, and Flow
}
\author{Educational Documentation}
\date{\today}

\begin{document}
\maketitle
\tableofcontents
\clearpage

%==============================================================================
\section{Introduction: What Is This Project?}
%==============================================================================

This project is a \textbf{backtesting platform} for trading strategies. Before we dive deep, let's understand what that means.

\begin{keyterm}{Backtesting}
\textbf{Backtesting} is the process of testing a trading strategy on historical data to see how it would have performed. Think of it as a time machine for your investment ideas—you can see what would have happened if you had used your strategy in the past.
\end{keyterm}

\subsection{Why Do We Need This?}

Imagine you have an idea: "I'll buy a stock when its price crosses above its 50-day average and sell when it crosses below." Before risking real money, you want to know:
\begin{itemize}
  \item Would this strategy have made money in the past?
  \item How much would I have gained or lost?
  \item How risky is this strategy?
  \item When does it perform well, and when does it fail?
\end{itemize}

This project provides:
\begin{enumerate}
  \item A \textbf{framework} to define strategies
  \item An \textbf{engine} to simulate trading
  \item A \textbf{web interface} to configure and run tests
  \item \textbf{Reports and visualizations} to understand results
\end{enumerate}

%==============================================================================
\section{Key Terms and Concepts}
%==============================================================================

Let's define all the important terms you'll encounter in this project.

\subsection{Core Trading Concepts}

\begin{keyterm}{Strategy}
A \textbf{strategy} is a set of rules that determines when to buy or sell an asset. For example:
\begin{itemize}
  \item Buy when RSI < 30 (oversold)
  \item Sell when RSI > 70 (overbought)
\end{itemize}
In code, a strategy is a Python class with methods like \texttt{generate\_signals(data)}.
\end{keyterm}

\begin{keyterm}{Portfolio}
A \textbf{portfolio} tracks your current holdings and cash over time. It includes:
\begin{itemize}
  \item Current cash balance
  \item Shares held
  \item Total portfolio value (cash + holdings)
  \item Historical performance
\end{itemize}
\end{keyterm}

\begin{keyterm}{Trade}
A \textbf{trade} is a single buy or sell transaction. Each trade has:
\begin{itemize}
  \item Date/time
  \item Type (BUY or SELL)
  \item Price
  \item Number of shares
  \item Commission paid
\end{itemize}
\end{keyterm}

\begin{keyterm}{Signal}
A \textbf{signal} is an indicator from your strategy. Common signals:
\begin{itemize}
  \item BUY: Enter a position
  \item SELL: Exit a position
  \item HOLD: Do nothing
\end{itemize}
\end{keyterm}

\subsection{Performance Metrics}

\begin{keyterm}{Total Return}
The percentage gain or loss over the entire backtest period.
\[
\text{Total Return} = \frac{\text{Final Value} - \text{Initial Capital}}{\text{Initial Capital}} \times 100\%
\]
\textbf{Example:} Start with \$100,000, end with \$118,200 → 18.2\% return
\end{keyterm}

\begin{keyterm}{Sharpe Ratio}
Measures risk-adjusted returns. Higher is better.
\[
\text{Sharpe Ratio} = \frac{\text{Average Return} - \text{Risk-Free Rate}}{\text{Standard Deviation of Returns}}
\]
\textbf{Interpretation:}
\begin{itemize}
  \item < 1: Poor
  \item 1-2: Good
  \item 2-3: Very good
  \item > 3: Excellent
\end{itemize}
\end{keyterm}

\begin{keyterm}{Maximum Drawdown}   
The largest peak-to-trough decline in portfolio value. Shows the worst loss you would have experienced.
\[
\text{Drawdown} = \frac{\text{Trough Value} - \text{Peak Value}}{\text{Peak Value}} \times 100\%
\]
\textbf{Example:} Your portfolio peaked at \$120,000 then dropped to \$109,600 → -8.7\% drawdown
\end{keyterm}

\begin{keyterm}{Win Rate}
The percentage of profitable trades.
\[
\text{Win Rate} = \frac{\text{Number of Winning Trades}}{\text{Total Trades}} \times 100\%
\]
\textbf{Example:} 66 wins out of 124 trades → 53.2\% win rate
\end{keyterm}

\begin{keyterm}{Sortino Ratio}
Similar to Sharpe but only penalizes downside volatility (losses). Better for strategies with asymmetric returns.
\end{keyterm}

\subsection{Trading Realism}

\begin{keyterm}{Commission}
The fee you pay to your broker for each trade. Typically 0.1\% to 0.5\% per trade.
\[
\text{Commission} = \text{Trade Value} \times \text{Commission Rate}
\]
\textbf{Example:} Buy \$10,000 worth of stock at 0.1\% commission → pay \$10
\end{keyterm}

\begin{keyterm}{Slippage}
The difference between the expected price and the actual execution price. Market orders often fill at slightly worse prices.
\[
\text{Execution Price} = \text{Signal Price} \times (1 + \text{Slippage Factor})
\]
\textbf{Example:} Expected to buy at \$100, but slippage of 0.05\% means you pay \$100.05
\end{keyterm}

\begin{keyterm}{Position Sizing}
How much of your portfolio to allocate to each trade. Common approaches:
\begin{itemize}
  \item Fixed percentage (e.g., 10\% per trade)
  \item Equal dollar amount
  \item Risk-based (e.g., 2\% risk per trade)
\end{itemize}
\end{keyterm}

\subsection{Technical Analysis}

\begin{keyterm}{Moving Average (MA)}
The average price over a specified number of periods. Smooths out price fluctuations.
\[
MA_n = \frac{1}{n}\sum_{i=0}^{n-1} P_{t-i}
\]
\textbf{Example:} 50-day MA is the average closing price over the last 50 days
\end{keyterm}

\begin{keyterm}{Moving Average Crossover}
A strategy where:
\begin{itemize}
  \item BUY when short-term MA crosses above long-term MA (golden cross)
  \item SELL when short-term MA crosses below long-term MA (death cross)
\end{itemize}
\end{keyterm}

\begin{keyterm}{RSI (Relative Strength Index)}
Momentum oscillator ranging from 0 to 100. Measures speed and magnitude of price changes.
\begin{itemize}
  \item RSI < 30: Oversold (potential buy)
  \item RSI > 70: Overbought (potential sell)
\end{itemize}
\end{keyterm}

%==============================================================================
\section{Project Architecture}
%==============================================================================

The project consists of three main components:

\subsection{System Components Diagram}

\begin{verbatim}
┌─────────────────────────────────────────────────────────────┐
│                      USER INTERFACE                          │
│  ┌──────────────┐  ┌──────────────┐  ┌──────────────┐      │
│  │  Dashboard   │  │  Strategy    │  │   Results    │      │
│  │    Screen    │  │Configuration │  │    Screen    │      │
│  └──────────────┘  └──────────────┘  └──────────────┘      │
│         React + Vite + TypeScript                           │
└─────────────────────────────────────────────────────────────┘
                           ↕ HTTP/JSON
┌─────────────────────────────────────────────────────────────┐
│                     BACKEND API                              │
│  ┌──────────────────────────────────────────────────────┐  │
│  │           FastAPI REST Endpoints                      │  │
│  │  /api/strategies, /api/backtests, etc.              │  │
│  └──────────────────────────────────────────────────────┘  │
│         FastAPI + Pydantic + BackgroundTasks                │
└─────────────────────────────────────────────────────────────┘
                           ↕ Python imports
┌─────────────────────────────────────────────────────────────┐
│                  BACKTESTING ENGINE                          │
│  ┌────────────┐  ┌────────────┐  ┌────────────┐           │
│  │    Data    │  │ Strategies │  │   Engine   │           │
│  │ Ingestion  │  │  (MA, RSI) │  │ Simulation │           │
│  └────────────┘  └────────────┘  └────────────┘           │
│  ┌────────────┐  ┌────────────┐                            │
│  │  Metrics   │  │ Reporting/ │                            │
│  │Calculator  │  │    Plots   │                            │
│  └────────────┘  └────────────┘                            │
│         Pandas + NumPy + Matplotlib                         │
└─────────────────────────────────────────────────────────────┘
\end{verbatim}

\subsection{Component Details}

\subsubsection{Frontend (React + Vite)}

\textbf{Purpose:} Provides a modern web interface for users to interact with the system.

\textbf{Key files:}
\begin{itemize}
  \item \texttt{src/components/Dashboard.tsx}: Shows KPIs and recent backtests
  \item \texttt{src/components/StrategyConfiguration.tsx}: Form to define strategies
  \item \texttt{src/components/BacktestExecution.tsx}: Start/monitor backtests
  \item \texttt{src/components/ReportScreen.tsx}: View results with charts
  \item \texttt{src/lib/api.ts}: API client for backend communication
\end{itemize}

\textbf{Technologies:}
\begin{itemize}
  \item React: UI framework
  \item Vite: Fast build tool
  \item TypeScript: Type-safe JavaScript
  \item Radix UI: Accessible component primitives
  \item shadcn/ui: Pre-built component patterns
  \item Recharts: Charting library
  \item Tailwind CSS: Utility-first styling
\end{itemize}

\subsubsection{Backend (FastAPI)}

\textbf{Purpose:} REST API that orchestrates backtests and serves results.

\textbf{Key file:} \texttt{ui/backend.py}

\textbf{Main responsibilities:}
\begin{enumerate}
  \item Accept strategy configurations
  \item Validate input data
  \item Run backtests in background tasks
  \item Track progress and status
  \item Serialize results to JSON
  \item Serve data to frontend
\end{enumerate}

\textbf{Technologies:}
\begin{itemize}
  \item FastAPI: Modern Python web framework
  \item Pydantic: Data validation
  \item Uvicorn: ASGI server
  \item BackgroundTasks: Async task execution
  \item CORS middleware: Cross-origin requests
\end{itemize}

\subsubsection{Backtesting Engine (Python)}

\textbf{Purpose:} Core logic for simulating trading strategies.

\textbf{Key modules:}
\begin{itemize}
  \item \texttt{backtesting\_engine/data/ingestion.py}: Load/generate price data
  \item \texttt{backtesting\_engine/strategies/}: Strategy implementations
  \item \texttt{backtesting\_engine/core/engine.py}: Main simulation engine
  \item \texttt{backtesting\_engine/analytics/metrics.py}: Performance calculations
  \item \texttt{backtesting\_engine/visualization/reporting.py}: Charts and reports
\end{itemize}

\textbf{Technologies:}
\begin{itemize}
  \item Pandas: Data manipulation
  \item NumPy: Numerical computations
  \item Matplotlib: Visualization
\end{itemize}

\subsubsection{Research Notebook}

\textbf{Purpose:} Interactive exploration and development.

\textbf{Key file:} \texttt{backtest.ipynb}

\textbf{Use cases:}
\begin{itemize}
  \item Prototype new strategies
  \item Analyze results in detail
  \item Create custom visualizations
  \item Document experiments
\end{itemize}

%==============================================================================
\section{Complete Project Flow}
%==============================================================================

Let's walk through the entire flow from strategy creation to results analysis.

\subsection{Flow Overview}

\begin{verbatim}
1. User defines strategy → 2. Frontend validates → 3. API receives
        ↓                            ↓                    ↓
4. Backend creates strategy instance ← ← ← ← ← ← ← ← ← ← ┘
        ↓
5. Load historical data
        ↓
6. Initialize portfolio (cash + holdings)
        ↓
7. FOR EACH trading day:
    a. Calculate indicators (MA, RSI, etc.)
    b. Generate signal (BUY/SELL/HOLD)
    c. Execute trades (with commission/slippage)
    d. Update portfolio
        ↓
8. Calculate performance metrics
        ↓
9. Generate visualizations
        ↓
10. Serialize results to JSON
        ↓
11. Frontend displays results
\end{verbatim}

\subsection{Detailed Flow Steps}

\begin{flowstep}{1: User Defines Strategy}
User fills out the strategy configuration form:
\begin{itemize}
  \item \textbf{Name:} "My MA Crossover v1"
  \item \textbf{Description:} "50/200 MA crossover on daily data"
  \item \textbf{Asset Class:} Stocks
  \item \textbf{Timeframe:} 1 Day
  \item \textbf{Entry Signal:} MA Crossover
  \item \textbf{Exit Signal:} MA Crossover (reverse)
  \item \textbf{Stop Loss:} 5\%
  \item \textbf{Take Profit:} 10\%
  \item \textbf{Position Size:} 10\% of portfolio
\end{itemize}
\end{flowstep}

\begin{flowstep}{2: Frontend Validation}
Client-side checks before submitting:
\begin{lstlisting}[style=pythoncode]
if (!formData.name.trim()) {
  toast({ title: "Error", 
          description: "Strategy name required" });
  return;
}
if (!formData.entry_signal) {
  toast({ title: "Error", 
          description: "Entry signal required" });
  return;
}
\end{lstlisting}
\end{flowstep}

\begin{flowstep}{3: API Request}
Frontend sends POST request to backend:
\begin{lstlisting}[style=jsoncode]
POST /api/strategies
{
  "name": "My MA Crossover v1",
  "description": "50/200 MA crossover on daily data",
  "asset_class": "stocks",
  "timeframe": "1d",
  "entry_signal": "ma-crossover",
  "exit_signal": "ma-crossover",
  "stop_loss": 5.0,
  "take_profit": 10.0,
  "position_size": 10.0,
  "custom_code": null
}
\end{lstlisting}
\end{flowstep}

\begin{flowstep}{4: Backend Creates Strategy}
Backend validates and instantiates strategy:
\begin{lstlisting}[style=pythoncode]
@app.post("/api/strategies")
async def create_strategy(strategy: StrategyConfig):
    engine = manager.get_engine()
    
    if strategy.entry_signal == "ma-crossover":
        new_strategy = MovingAverageCrossover(
            short_window=50, 
            long_window=200
        )
    elif strategy.entry_signal == "rsi-oversold":
        new_strategy = RSIStrategy(
            period=14, 
            oversold=30, 
            overbought=70
        )
    
    engine.add_strategy(new_strategy)
    return {"message": "Strategy created successfully"}
\end{lstlisting}
\end{flowstep}

\begin{flowstep}{5: Load Historical Data}
Data ingestion module loads or generates price data:
\begin{lstlisting}[style=pythoncode]
data_ingestion = DataIngestion()
data = data_ingestion.generate_sample_data(
    start_date="2024-01-01",
    end_date="2024-12-31"
)
# Returns DataFrame with columns:
# Date, Open, High, Low, Close, Volume
\end{lstlisting}
\end{flowstep}

\begin{flowstep}{6: Initialize Portfolio}
Engine sets up initial portfolio state:
\begin{lstlisting}[style=pythoncode]
self.initial_capital = 100000.0
self.cash = 100000.0
self.holdings = 0
self.portfolio_value = 100000.0
self.trades = []
self.portfolio_history = []
\end{lstlisting}
\end{flowstep}

\begin{flowstep}{7: Main Simulation Loop}
For each day in the dataset:

\textbf{Step 7a: Calculate Indicators}
\begin{lstlisting}[style=pythoncode]
# Moving Average Strategy
data['MA_short'] = data['Close'].rolling(
    window=short_window
).mean()
data['MA_long'] = data['Close'].rolling(
    window=long_window
).mean()
\end{lstlisting}

\textbf{Step 7b: Generate Signals}
\begin{lstlisting}[style=pythoncode]
def generate_signals(self, data):
    signals = pd.DataFrame(index=data.index)
    signals['signal'] = 0
    
    # Buy when short MA crosses above long MA
    signals['signal'] = np.where(
        data['MA_short'] > data['MA_long'], 1, 0
    )
    
    # Generate trading orders
    signals['positions'] = signals['signal'].diff()
    return signals
\end{lstlisting}

\textbf{Step 7c: Execute Trades}
\begin{lstlisting}[style=pythoncode]
if signal == 1:  # BUY
    # Calculate position size
    position_value = self.cash * position_size
    shares = position_value / execution_price
    
    # Apply slippage
    execution_price *= (1 + slippage_factor)
    
    # Calculate commission
    commission = position_value * commission_rate
    
    # Update portfolio
    self.cash -= (position_value + commission)
    self.holdings += shares
    
    # Record trade
    self.trades.append({
        'date': current_date,
        'type': 'BUY',
        'price': execution_price,
        'shares': shares,
        'commission': commission
    })
\end{lstlisting}

\textbf{Step 7d: Update Portfolio}
\begin{lstlisting}[style=pythoncode]
# Calculate current portfolio value
holdings_value = self.holdings * current_price
portfolio_value = self.cash + holdings_value

# Record in history
self.portfolio_history.append({
    'date': current_date,
    'cash': self.cash,
    'holdings': self.holdings,
    'holdings_value': holdings_value,
    'total_value': portfolio_value
})
\end{lstlisting}
\end{flowstep}

\begin{flowstep}{8: Calculate Metrics}
After simulation completes, calculate performance:
\begin{lstlisting}[style=pythoncode]
# Total Return
total_return = (
    (final_value - initial_capital) / initial_capital
) * 100

# Sharpe Ratio
returns = portfolio['daily_return']
sharpe = (
    returns.mean() / returns.std()
) * np.sqrt(252)  # Annualized

# Maximum Drawdown
cumulative = (1 + returns).cumprod()
peak = cumulative.expanding(min_periods=1).max()
drawdown = (cumulative - peak) / peak
max_drawdown = drawdown.min() * 100

# Win Rate
winning_trades = sum(1 for t in trades 
                     if t['profit'] > 0)
win_rate = (winning_trades / len(trades)) * 100
\end{lstlisting}
\end{flowstep}

\begin{flowstep}{9: Generate Visualizations}
Create charts for analysis:
\begin{lstlisting}[style=pythoncode]
# Equity curve
plt.plot(portfolio['date'], 
         portfolio['total_value'])
plt.title('Portfolio Equity Curve')
plt.xlabel('Date')
plt.ylabel('Portfolio Value ($)')

# Returns distribution
plt.hist(portfolio['daily_return'], bins=50)
plt.title('Daily Returns Distribution')
\end{lstlisting}
\end{flowstep}

\begin{flowstep}{10: Serialize Results}
Convert results to JSON for API response:
\begin{lstlisting}[style=pythoncode]
results = {
    "backtest_id": backtest_id,
    "strategy_name": strategy_name,
    "metrics": {
        "total_return": float(total_return),
        "sharpe_ratio": float(sharpe_ratio),
        "max_drawdown": float(max_drawdown),
        "win_rate": float(win_rate)
    },
    "portfolio_data": [
        {
            "date": row['date'].strftime("%Y-%m-%d"),
            "portfolio_value": float(row['total_value'])
        }
        for _, row in portfolio.iterrows()
    ],
    "trades": [/* ... */]
}
\end{lstlisting}
\end{flowstep}

\begin{flowstep}{11: Display Results}
Frontend fetches and displays:
\begin{lstlisting}[style=pythoncode]
const results = await apiClient.getBacktestResults(
    backtestId
);

// Render KPI cards
<div>Total Return: {results.metrics.total_return}%</div>
<div>Sharpe Ratio: {results.metrics.sharpe_ratio}</div>

// Render chart
<LineChart data={results.portfolio_data}>
  <Line dataKey="portfolio_value" />
</LineChart>
\end{lstlisting}
\end{flowstep}

%==============================================================================
\section{Code Structure Deep Dive}
%==============================================================================

\subsection{Module Organization}

\begin{verbatim}
BacktestingEngine/
├── backtesting_engine/         # Core engine
│   ├── __init__.py
│   ├── data/
│   │   ├── __init__.py
│   │   └── ingestion.py        # Data loading
│   ├── strategies/
│   │   ├── __init__.py
│   │   ├── base.py             # Base strategy class
│   │   ├── moving_average.py   # MA strategy
│   │   └── rsi.py              # RSI strategy
│   ├── core/
│   │   ├── __init__.py
│   │   └── engine.py           # Main engine
│   ├── analytics/
│   │   ├── __init__.py
│   │   └── metrics.py          # Performance metrics
│   └── visualization/
│       ├── __init__.py
│       └── reporting.py        # Charts & reports
├── ui/                         # Web interface
│   ├── src/
│   │   ├── components/         # React components
│   │   ├── lib/                # Utilities & API
│   │   └── main.tsx            # Entry point
│   ├── backend.py              # FastAPI server
│   ├── package.json
│   └── vite.config.ts
├── tests/                      # Unit tests
├── docs/                       # Documentation
├── backtest.ipynb              # Jupyter notebook
└── README.md
\end{verbatim}

\subsection{Key Classes and Their Roles}

\subsubsection{DataIngestion}

\textbf{Purpose:} Load or generate price data for backtesting.

\textbf{Methods:}
\begin{lstlisting}[style=pythoncode]
class DataIngestion:
    def load_from_csv(self, filepath):
        """Load data from CSV file"""
        
    def generate_sample_data(self, start_date, end_date):
        """Generate synthetic OHLCV data"""
        
    def fetch_from_api(self, symbol, start, end):
        """Fetch from external API (future)"""
\end{lstlisting}

\subsubsection{Strategy (Base Class)}

\textbf{Purpose:} Abstract base for all strategies.

\textbf{Interface:}
\begin{lstlisting}[style=pythoncode]
class Strategy(ABC):
    @abstractmethod
    def generate_signals(self, data):
        """
        Generate trading signals from data
        Returns: DataFrame with 'signal' column
        """
        pass
    
    @abstractmethod
    def calculate_indicators(self, data):
        """
        Add technical indicators to data
        """
        pass
\end{lstlisting}

\subsubsection{MovingAverageCrossover}

\textbf{Purpose:} Implements MA crossover strategy.

\textbf{Parameters:}
\begin{itemize}
  \item \texttt{short\_window}: Fast MA period (e.g., 50)
  \item \texttt{long\_window}: Slow MA period (e.g., 200)
\end{itemize}

\textbf{Logic:}
\begin{lstlisting}[style=pythoncode]
class MovingAverageCrossover(Strategy):
    def __init__(self, short_window=50, long_window=200):
        self.short_window = short_window
        self.long_window = long_window
        
    def generate_signals(self, data):
        # Calculate MAs
        data['MA_short'] = data['Close'].rolling(
            self.short_window
        ).mean()
        data['MA_long'] = data['Close'].rolling(
            self.long_window
        ).mean()
        
        # Generate signals
        signals = pd.DataFrame(index=data.index)
        signals['signal'] = np.where(
            data['MA_short'] > data['MA_long'], 1, 0
        )
        signals['positions'] = signals['signal'].diff()
        
        return signals
\end{lstlisting}

\subsubsection{BacktestingEngine}

\textbf{Purpose:} Main simulation engine that executes strategies.

\textbf{Key methods:}
\begin{lstlisting}[style=pythoncode]
class BacktestingEngine:
    def __init__(self, initial_capital=100000):
        self.initial_capital = initial_capital
        self.strategies = {}
        
    def add_strategy(self, strategy):
        """Register a strategy"""
        
    def run_backtest(self, data, strategy_name, 
                    commission=0.001, slippage=0.0005):
        """
        Execute backtest for a strategy
        Returns: dict with portfolio, trades, metrics
        """
        
    def _execute_trade(self, signal, price, date):
        """Execute a single trade with realistic costs"""
        
    def _update_portfolio(self, current_price):
        """Update portfolio value and history"""
\end{lstlisting}

\subsubsection{Reporting}

\textbf{Purpose:} Calculate metrics and generate visualizations.

\textbf{Metrics calculated:}
\begin{lstlisting}[style=pythoncode]
class Reporting:
    @staticmethod
    def calculate_metrics(portfolio, trades):
        return {
            'total_return_pct': ...,
            'annualized_return_pct': ...,
            'sharpe_ratio': ...,
            'sortino_ratio': ...,
            'max_drawdown_pct': ...,
            'win_rate_pct': ...,
            'total_trades': len(trades),
            'avg_trade_return': ...,
            'profit_factor': ...
        }
    
    @staticmethod
    def plot_equity_curve(portfolio, metrics):
        """Generate equity curve visualization"""
        
    @staticmethod
    def plot_returns_distribution(portfolio):
        """Plot returns histogram and statistics"""
\end{lstlisting}

%==============================================================================
\section{API Reference}
%==============================================================================

\subsection{Endpoint Categories}

\subsubsection{Dashboard Endpoints}

\textbf{GET /api/dashboard/stats}

Returns summary statistics.

\textit{Response:}
\begin{lstlisting}[style=jsoncode]
{
  "total_strategies": 12,
  "active_backtests": 0,
  "avg_sharpe_ratio": 1.42,
  "best_strategy_roi": 23.4
}
\end{lstlisting}

\textbf{GET /api/dashboard/recent-backtests}

Returns recent backtest runs.

\textit{Response:}
\begin{lstlisting}[style=jsoncode]
[
  {
    "id": "BT-001",
    "strategy": "Mean Reversion v2",
    "status": "Completed",
    "sharpe": "1.34",
    "roi": "18.2%"
  }
]
\end{lstlisting}

\subsubsection{Strategy Endpoints}

\textbf{GET /api/strategies}

List available strategies.

\textbf{POST /api/strategies}

Create a new strategy.

\textit{Request body:}
\begin{lstlisting}[style=jsoncode]
{
  "name": "My Strategy",
  "description": "...",
  "asset_class": "stocks",
  "timeframe": "1d",
  "entry_signal": "ma-crossover",
  "exit_signal": "stop-loss",
  "stop_loss": 5.0,
  "take_profit": 10.0,
  "position_size": 10.0
}
\end{lstlisting}

\subsubsection{Backtest Endpoints}

\textbf{POST /api/backtests/start}

Start a new backtest.

\textit{Request:}
\begin{lstlisting}[style=jsoncode]
{
  "strategy_name": "ma_crossover",
  "start_date": "2024-01-01",
  "end_date": "2024-12-31",
  "initial_capital": 100000.0,
  "commission": 0.001,
  "slippage": 0.0005
}
\end{lstlisting}

\textit{Response:}
\begin{lstlisting}[style=jsoncode]
{
  "backtest_id": "abc123",
  "status": "started"
}
\end{lstlisting}

\textbf{GET /api/backtests/\{id\}/status}

Check backtest progress.

\textit{Response:}
\begin{lstlisting}[style=jsoncode]
{
  "id": "abc123",
  "status": "running",
  "progress": 67.5,
  "current_step": "Running simulation...",
  "estimated_time_remaining": "1m 23s"
}
\end{lstlisting}

\textbf{GET /api/backtests/\{id\}/results}

Retrieve completed backtest results.

\textit{Response:} See Section~\ref{sec:results-format}

\textbf{POST /api/backtests/\{id\}/pause}

Pause a running backtest.

\textbf{POST /api/backtests/\{id\}/stop}

Stop a running backtest.

\subsection{Results Format}
\label{sec:results-format}

Complete backtest results structure:

\begin{lstlisting}[style=jsoncode]
{
  "backtest_id": "abc123",
  "strategy_name": "MovingAverageCrossover",
  "metrics": {
    "total_return": 18.2,
    "annualized_return": 12.5,
    "sharpe_ratio": 1.34,
    "sortino_ratio": 1.90,
    "max_drawdown": -8.7,
    "win_rate": 53.2,
    "total_trades": 124
  },
  "portfolio_data": [
    {
      "date": "2024-01-02",
      "portfolio_value": 100523.4,
      "daily_return": 0.005234,
      "cumulative_return": 0.005234
    },
    // ... more dates
  ],
  "trades": [
    {
      "id": "T001",
      "date": "2024-02-10",
      "type": "BUY",
      "price": 152.30,
      "shares": 65,
      "commission": 9.90
    },
    // ... more trades
  ],
  "completed_at": "2024-01-15T12:34:56Z"
}
\end{lstlisting}

%==============================================================================
\section{Running the Project}
%==============================================================================

\subsection{Prerequisites}

\begin{itemize}
  \item Python 3.8+
  \item Node.js 16+
  \item pip (Python package manager)
  \item npm (Node package manager)
\end{itemize}

\subsection{Installation Steps}

\subsubsection{Backend Setup}

\begin{lstlisting}[style=shellcode]
# 1. Create virtual environment
cd BacktestingEngine
python -m venv .venv

# 2. Activate virtual environment
source .venv/bin/activate  # Linux/Mac
# or
.venv\Scripts\activate     # Windows

# 3. Install Python dependencies
pip install -r requirements.txt

# 4. Start the backend server
cd ui
uvicorn backend:app --reload --host 0.0.0.0 --port 8000
\end{lstlisting}

Backend will be available at: \url{http://localhost:8000}

\subsubsection{Frontend Setup}

\begin{lstlisting}[style=shellcode]
# 1. Navigate to UI directory
cd ui

# 2. Install Node dependencies
npm install

# 3. Start development server
npm run dev
\end{lstlisting}

Frontend will be available at: \url{http://localhost:3000}

\subsubsection{Notebook Setup}

\begin{lstlisting}[style=shellcode]
# 1. Ensure virtual environment is activated
source .venv/bin/activate

# 2. Install Jupyter (if not already)
pip install jupyter

# 3. Open notebook
jupyter notebook backtest.ipynb
\end{lstlisting}

\subsection{Testing the Setup}

\textbf{Test backend:}
\begin{lstlisting}[style=shellcode]
curl http://localhost:8000/
# Should return: {"message":"Backtesting Engine API",...}
\end{lstlisting}

\textbf{Test strategies endpoint:}
\begin{lstlisting}[style=shellcode]
curl http://localhost:8000/api/strategies
# Should return array of available strategies
\end{lstlisting}

%==============================================================================
\section{Example: Complete Backtest Workflow}
%==============================================================================

Let's walk through a complete example from start to finish.

\subsection{Scenario}

\textbf{Goal:} Test a 50/200 MA crossover strategy on synthetic stock data from Jan 1, 2024 to Dec 31, 2024.

\subsection{Step-by-Step Walkthrough}

\begin{example}{Step 1: Define Strategy in UI}
Open \url{http://localhost:3000} and fill in the form:
\begin{itemize}
  \item Name: "Test MA Crossover"
  \item Description: "50/200 day MA crossover"
  \item Asset Class: Stocks
  \item Timeframe: 1 Day
  \item Entry Signal: MA Crossover
  \item Exit Signal: MA Crossover
  \item Stop Loss: 5\%
  \item Position Size: 10\%
\end{itemize}
Click "Save Strategy"
\end{example}

\begin{example}{Step 2: Configure Backtest}
Navigate to "Backtest Execution" screen:
\begin{itemize}
  \item Select Strategy: "Test MA Crossover"
  \item Start Date: 2024-01-01
  \item End Date: 2024-12-31
  \item Initial Capital: \$100,000
  \item Commission: 0.1\%
  \item Slippage: 0.05\%
\end{itemize}
Click "Start Backtest"
\end{example}

\begin{example}{Step 3: Monitor Progress}
Watch the progress bar:
\begin{itemize}
  \item 10\%: Loading data...
  \item 30\%: Initializing strategy...
  \item 50\%: Running simulation...
  \item 80\%: Calculating metrics...
  \item 100\%: Backtest completed successfully
\end{itemize}
\end{example}

\begin{example}{Step 4: View Results}
Results screen shows:

\textbf{Key Metrics:}
\begin{itemize}
  \item Total Return: 18.2\%
  \item Sharpe Ratio: 1.34
  \item Max Drawdown: -8.7\%
  \item Win Rate: 53.2\%
  \item Total Trades: 124
\end{itemize}

\textbf{Charts:}
\begin{itemize}
  \item Equity curve: Shows steady growth
  \item Drawdown: Largest decline was -8.7\% in August
  \item Returns distribution: Normally distributed
\end{itemize}

\textbf{Recent Trades:}
\begin{verbatim}
T124 | 2024-12-20 | SELL | $158.30 | 63 shares
T123 | 2024-11-15 | BUY  | $147.80 | 67 shares
T122 | 2024-10-22 | SELL | $152.10 | 65 shares
\end{verbatim}
\end{example}

\begin{example}{Step 5: Analysis}
\textbf{Interpretation:}
\begin{itemize}
  \item Strategy is profitable (18.2\% return)
  \item Risk-adjusted performance is good (Sharpe 1.34)
  \item Drawdowns are manageable (-8.7\% max)
  \item Win rate slightly above 50\% (53.2\%)
\end{itemize}

\textbf{Next steps:}
\begin{itemize}
  \item Test on real historical data
  \item Optimize MA periods (30/100, 20/50, etc.)
  \item Add stop-loss/take-profit rules
  \item Test on different asset classes
\end{itemize}
\end{example}

%==============================================================================
\section{Advanced Topics}
%==============================================================================

\subsection{Adding a Custom Strategy}

\textbf{Step 1:} Create new strategy file

\begin{lstlisting}[style=pythoncode]
# backtesting_engine/strategies/bollinger_bands.py

from .base import Strategy
import pandas as pd
import numpy as np

class BollingerBands(Strategy):
    def __init__(self, period=20, std_dev=2):
        self.period = period
        self.std_dev = std_dev
        self.name = "BollingerBands"
        
    def calculate_indicators(self, data):
        # Calculate Bollinger Bands
        data['BB_middle'] = data['Close'].rolling(
            self.period
        ).mean()
        std = data['Close'].rolling(self.period).std()
        data['BB_upper'] = data['BB_middle'] + (
            self.std_dev * std
        )
        data['BB_lower'] = data['BB_middle'] - (
            self.std_dev * std
        )
        return data
        
    def generate_signals(self, data):
        data = self.calculate_indicators(data)
        signals = pd.DataFrame(index=data.index)
        
        # Buy when price touches lower band
        # Sell when price touches upper band
        signals['signal'] = 0
        signals.loc[
            data['Close'] <= data['BB_lower'], 'signal'
        ] = 1
        signals.loc[
            data['Close'] >= data['BB_upper'], 'signal'
        ] = -1
        
        signals['positions'] = signals['signal'].diff()
        return signals
\end{lstlisting}

\textbf{Step 2:} Register in backend

\begin{lstlisting}[style=pythoncode]
# ui/backend.py

from backtesting_engine import BollingerBands

@app.post("/api/strategies")
async def create_strategy(strategy: StrategyConfig):
    # ... existing code ...
    elif strategy.entry_signal == "bollinger-bands":
        new_strategy = BollingerBands(
            period=20, 
            std_dev=2
        )
    # ... rest of code ...
\end{lstlisting}

\subsection{Performance Optimization Tips}

\begin{warning}
For large datasets (>10 years daily data):
\begin{enumerate}
  \item Use vectorized operations (NumPy/Pandas)
  \item Avoid Python loops where possible
  \item Cache indicator calculations
  \item Consider using Numba for hot paths
  \item Batch process multiple strategies in parallel
\end{enumerate}
\end{warning}

\subsection{Common Pitfalls and Solutions}

\begin{enumerate}
  \item \textbf{Look-ahead bias}: Using future data in decisions
  \begin{itemize}
    \item Solution: Ensure indicators only use past data
  \end{itemize}
  
  \item \textbf{Overfitting}: Strategy works on backtest but fails live
  \begin{itemize}
    \item Solution: Use out-of-sample testing, walk-forward analysis
  \end{itemize}
  
  \item \textbf{Ignoring costs}: Unrealistic performance from zero costs
  \begin{itemize}
    \item Solution: Always include commission and slippage
  \end{itemize}
  
  \item \textbf{Survivorship bias}: Testing only on surviving companies
  \begin{itemize}
    \item Solution: Include delisted stocks in dataset
  \end{itemize}
\end{enumerate}

%==============================================================================
\section{Study Checklist}
%==============================================================================

Use this checklist to verify your understanding:

\subsection{Conceptual Understanding}
\begin{itemize}
  \item[$\square$] I understand what backtesting is and why it's important
  \item[$\square$] I can explain the difference between strategies
  \item[$\square$] I understand key metrics (Sharpe, drawdown, win rate)
  \item[$\square$] I know why commission and slippage matter
  \item[$\square$] I can interpret backtest results
\end{itemize}

\subsection{System Architecture}
\begin{itemize}
  \item[$\square$] I understand the three-tier architecture
  \item[$\square$] I know what each component does (frontend/backend/engine)
  \item[$\square$] I can trace a request from UI to results
  \item[$\square$] I understand the API endpoints
  \item[$\square$] I know where strategies are defined
\end{itemize}

\subsection{Implementation}
\begin{itemize}
  \item[$\square$] I can run the backend server
  \item[$\square$] I can run the frontend
  \item[$\square$] I can create a strategy through the UI
  \item[$\square$] I can run a backtest
  \item[$\square$] I can interpret the results
\end{itemize}

\subsection{Advanced}
\begin{itemize}
  \item[$\square$] I can add a custom strategy
  \item[$\square$] I understand the simulation loop
  \item[$\square$] I can modify metrics calculations
  \item[$\square$] I can add new visualizations
  \item[$\square$] I understand common pitfalls
\end{itemize}

%==============================================================================
\section{Further Resources}
%==============================================================================

\subsection{Recommended Reading}
\begin{itemize}
  \item \textit{Algorithmic Trading} by Ernest P. Chan
  \item \textit{Quantitative Trading} by Ernest P. Chan
  \item \textit{Python for Finance} by Yves Hilpisch
  \item FastAPI Documentation: \url{https://fastapi.tiangolo.com}
  \item React Documentation: \url{https://react.dev}
\end{itemize}

\subsection{Online Courses}
\begin{itemize}
  \item Coursera: "Machine Learning for Trading"
  \item Udemy: "Algorithmic Trading \& Quantitative Analysis"
  \item YouTube: "sentdex" Python finance tutorials
\end{itemize}

\subsection{Related Projects}
\begin{itemize}
  \item Zipline: Pythonic algorithmic trading library
  \item Backtrader: Feature-rich backtesting framework
  \item QuantConnect: Cloud-based backtesting platform
  \item PyAlgoTrade: Event-driven backtesting library
\end{itemize}

%==============================================================================
\section*{Compilation Instructions}
%==============================================================================

To compile this document:

\begin{lstlisting}[style=shellcode]
# Option 1: Using latexmk (recommended)
latexmk -pdf docs/study-guide.tex

# Option 2: Using pdflatex (run twice for ToC)
pdflatex docs/study-guide.tex
pdflatex docs/study-guide.tex

# Clean auxiliary files
latexmk -c docs/study-guide.tex
\end{lstlisting}

Output: \texttt{study-guide.pdf}

\end{document}
